\documentclass[11pt]{article}
\usepackage[margin=1in]{geometry}
\usepackage{amsmath,amssymb,amsthm,mathtools}
\usepackage{graphicx}
\usepackage{microtype}
\usepackage[hidelinks]{hyperref}

\newtheorem{theorem}{Theorem}
\newcommand{\C}{\mathbb C}
\newcommand{\tr}{\operatorname{Tr}}
\newcommand{\id}{\mathrm{Id}}

\title{An all-cuts PPT-entangled tripartite Werner state}
\author{Autonomous proof candidate for specialist verification}
\date{9 August 2026}

\begin{document}
\maketitle

\begin{abstract}
We answer affirmatively the open question in Remark~5.5 of
Park--Jung--Park--Youn: for every $d\ge3$, there is a
$U\otimes U\otimes U$-invariant state on $(\C^d)^{\otimes3}$ which is
positive under partial transpose and entangled across each of the three
bipartitions.  In the coordinates of Eggeling--Werner, the example is the
rational point
\[
 (r_+,r_-,r_1,r_2,r_3)=\left(\frac34,\frac1{32},0,0,\frac3{16}\right).
\]
It is invariant under the cyclic permutation of the three systems.  The
one-cut PPT criteria hold strictly, whereas the corresponding biseparability
criterion fails by $45/1024$.  A bounded search found no later explicit
resolution, so novelty remains subject to specialist review.
\end{abstract}

\section{Source question}

Park, Jung, Park, and Youn~\cite{PJYP} construct tripartite Werner states
which are PPT-entangled across two cuts but not PPT across the third.  In
Remark~5.5 on PDF page 25 they ask for a state which is PPT-entangled across
all three cuts.  The complete source passage is reproduced below.

\begin{center}
 \includegraphics[width=.94\textwidth]{figures/open_problem_crop.png}
\end{center}

In the source's terminology, a tripartite Werner state is a density operator
$\rho$ on $(\C^d)^{\otimes3}$ satisfying
\[
 (U\otimes U\otimes U)\rho(U\otimes U\otimes U)^*=\rho
 \qquad(U\in\mathcal U_d).
\]
It is, for example, $A$--$BC$ PPT if $\rho^{T_A}\ge0$, and it is
$A$--$BC$ entangled if it is not separable across that bipartition.

\section{Eggeling--Werner coordinates}

We recall only the exact criteria used below from Eggeling--Werner~\cite{EW}.
For $\sigma\in S_3$, let $V_\sigma$ permute the three tensor factors, and put
\begin{align*}
R_+&=\tfrac16(\id+V_{(12)}+V_{(23)}+V_{(13)}+V_{(123)}+V_{(132)}),\\
R_-&=\tfrac16(\id-V_{(12)}-V_{(23)}-V_{(13)}+V_{(123)}+V_{(132)}),\\
R_0&=\tfrac13(2\id-V_{(123)}-V_{(132)}),\\
R_1&=\tfrac13(2V_{(23)}-V_{(13)}-V_{(12)}),\\
R_2&=\tfrac1{\sqrt3}(V_{(12)}-V_{(13)}),\qquad
R_3=\tfrac{i}{\sqrt3}(V_{(123)}-V_{(132)}).
\end{align*}
For a Werner state define $r_k=\tr(\rho R_k)$.  Such a state is characterized
by
\[
 r_+,r_-,r_0\ge0,\qquad r_++r_-+r_0=1,\qquad
 r_1^2+r_2^2+r_3^2\le r_0^2.                                          \tag{1}
\]

For the cut $A$--$BC$, Lemma~10 of~\cite{EW} says that PPT is equivalent to
the state conditions together with
\begin{align}
0&\le r_-,&
0&\le1+r_1-r_+-r_-,\notag\\
0&\le1-r_1-5r_--r_+,&
0&\le-1-r_1+r_-+5r_+,                                                 \tag{2}\\
r_2^2+r_3^2&\le
 \frac{(1-r_1-5r_--r_+)(-1-r_1+r_-+5r_+)}3,\notag\\
r_2^2+r_3^2&\le(1-r_1-r_--r_+)(1+r_1-r_--r_+).\notag
\end{align}

We also need the negative branch of the biseparability criterion in
Theorem~5 of~\cite{EW}.  If $0\le r_-<1/3$ and
\[
 -1<u:=\frac{1+r_1-r_--2r_+}{1-3r_-}\le0,                             \tag{3}
\]
then $A$--$BC$ separability requires
\[
 3r_2^2+3r_3^2+(1+2r_1+r_--r_+)^2
 \le(2+r_1-4r_--2r_+)^2.                                             \tag{4}
\]
Together with the other branch and three isolated extreme points, the cited
theorem makes these conditions necessary and sufficient.  Our point lies
strictly inside the parameter range (3), so only (4) is relevant.

\section{Proof intuition}

The obstruction quoted in the source applies when a transposition symmetry is
present.  We instead impose only cyclic symmetry.  In the Eggeling--Werner
Bloch coordinates, setting $r_1=r_2=0$ makes the state invariant under the
three-cycle, so the three bipartitions have identical PPT and separability
status.  A nonzero $r_3$ breaks the transposition symmetries.  It remains only
to choose rational $r_+,r_-,r_3$ inside the PPT body but outside the
biseparable body; the point below does so with strict rational margins.

\section{Construction and proof}

\begin{theorem}\label{thm:main}
For every integer $d\ge3$, define
\begin{align}
\rho_d={}&\frac{9}{2d(d+1)(d+2)}R_+
 +\frac{3}{16d(d-1)(d-2)}R_- \notag\\
&+\frac{21}{64d(d^2-1)}R_0
 +\frac{9}{32d(d^2-1)}R_3.                                           \tag{5}
\end{align}
Then $\rho_d$ is a tripartite Werner state which is PPT and entangled across
each of $A$--$BC$, $B$--$AC$, and $C$--$AB$.
\end{theorem}

\begin{proof}
The trace dimensions of the three blocks give
\[
 (r_+,r_-,r_0,r_1,r_2,r_3)
 =\left(\frac34,\frac1{32},\frac7{32},0,0,\frac3{16}\right).           \tag{6}
\]
Indeed, the coefficient-to-expectation relations are
\[
 r_+=\frac{d(d+1)(d+2)}6c_+,\quad
 r_-=\frac{d(d-1)(d-2)}6c_-,\quad
 r_i=\frac{2d(d^2-1)}3c_i\quad(0\le i\le3).
\]
Thus $\tr\rho_d=1$.  Condition (1) holds because
\[
 r_3^2=\frac{36}{1024}<\frac{49}{1024}=r_0^2,                          \tag{7}
\]
so $\rho_d\ge0$.  Formula (5) is a linear combination of permutation
operators and therefore is $U\otimes U\otimes U$ invariant.

We first examine the cut $A$--$BC$.  At (6), the nontrivial linear quantities
in (2) are
\[
 1-r_+-r_-=\frac7{32},\qquad
 1-5r_--r_+=\frac3{32},\qquad
 -1+r_-+5r_+=\frac{89}{32}.                                          \tag{8}
\]
Moreover,
\[
 r_2^2+r_3^2=\frac{36}{1024},\qquad
 \frac{(3/32)(89/32)}3=\frac{89}{1024},\qquad
 \left(\frac7{32}\right)^2=\frac{49}{1024}.                          \tag{9}
\]
All inequalities in (2) therefore hold, and $\rho_d$ is $A$--$BC$ PPT.

It is not $A$--$BC$ separable.  The branch parameter (3) equals
\[
 u=\frac{1-1/32-2(3/4)}{1-3/32}=-\frac{17}{29},                        \tag{10}
\]
which lies strictly between $-1$ and $0$.  But the two sides of (4) are
respectively
\[
 3\left(\frac3{16}\right)^2
 +\left(1+\frac1{32}-\frac34\right)^2
 =\frac{189}{1024},\qquad
 \left(2-\frac4{32}-\frac32\right)^2
 =\frac{144}{1024}.                                                    \tag{11}
\]
The necessary inequality fails strictly, by $45/1024$.  Hence $\rho_d$ is
$A$--$BC$ entangled.

Finally, conjugation by $V_{(123)}$ fixes $R_+,R_-,R_0$, and $R_3$, hence
fixes $\rho_d$.  This conjugation cyclically permutes the three tensor factors,
and therefore carries the PPT cone and the separable cone for one cut onto
those for the next cut.  The $A$--$BC$ conclusions consequently hold for
$B$--$AC$ and $C$--$AB$ as well.
\end{proof}

\section{Why the quoted obstruction does not apply}

Remark~5.5 of~\cite{PJYP} rules out an all-cuts example under any of three
coefficient symmetries, each of which also requires
$a_{(123)}=a_{(132)}$.  In (5), the $R_3$ term contributes
\[
 \frac{9i}{32\sqrt3\,d(d^2-1)}
 \bigl(V_{(123)}-V_{(132)}\bigr),
\]
so these two coefficients are unequal.  This is precisely the
transposition-symmetry-breaking direction invisible in the ruled-out slices.

\section{Verification and novelty limitations}

The included script \texttt{code/verify\_candidate.py} checks (7)--(11) using
exact rational arithmetic.  It also constructs (5) directly for $d=3,4,5$
and checks trace, positivity, coordinates, cyclic invariance, and all three
partial transposes numerically.  These matrix checks are safeguards; the proof
above is dimension-independent and analytic.

A bounded exact-phrase and keyword search found the source question,
Eggeling--Werner's characterization, generic all-cuts PPT-entangled states,
and broader tripartite-entanglement work, but no later explicit solution in the
tripartite Werner class.  This is not an exhaustive bibliographic proof.  A
specialist should check forward citations and databases before any public
novelty claim.

\begin{thebibliography}{9}
\bibitem{PJYP}
S.-J. Park, Y.-G. Jung, J. Park, and S.-G. Youn,
\emph{A universal framework for entanglement detection under group symmetry},
arXiv:2301.03849v2 (2023), Remark 5.5.

\bibitem{EW}
T. Eggeling and R. F. Werner,
\emph{Separability properties of tripartite states with
$U\otimes U\otimes U$ symmetry},
Phys. Rev. A 63 (2001), 042111; arXiv:quant-ph/0010096.
\end{thebibliography}

\end{document}
